\documentclass[11pt]{article}
\usepackage{amsmath,amssymb,amsthm}
\newtheorem{theorem}{Theorem}
\newtheorem{proposition}[theorem]{Proposition}
\title{Asymmetric Anchors and Their Readout Dimension}
\author{}
\date{November 24, 2026}
\begin{document}
\maketitle

\paragraph{Status.}
These are complete written finite statements, not Lean declarations.
They audit a proposed compact reciprocal decoder. They neither change
the attention model nor establish an asymptotic head bound.

\section{An arbitrary assignment height from asymmetric zero anchors}

Fix distinct poles $\lambda_1,\ldots,\lambda_H$ and distinct anchor
arguments $d_1,\ldots,d_H$, all strictly above every pole. Fix a further
argument $D$ strictly above every pole and distinct from all anchors.
For every permutation $\pi$, consider the unique residue vector $c(\pi)$
solving
\[
 \sum_{h=1}^H\frac{a_{ih}+c_h(\pi)}{d_i-\lambda_{\pi(h)}}=0,
 \qquad 1\leq i\leq H. \tag{1}
\]
Here the anchor numerators $a_{ih}$ are fixed independently of $\pi$.

\begin{theorem}[Exact assignment-height realization]
For every real cost matrix $K=(K_{hj})$, the anchor numerators can be
chosen so that the unique solutions of (1) satisfy
\[
 \sum_h\frac{c_h(\pi)}{D-\lambda_{\pi(h)}}
 =-\sum_h K_{h,\pi(h)}
 \qquad\hbox{for every permutation }\pi. \tag{2}
\]
This statement remains true when all poles and all arguments are confined
to arbitrarily small, disjoint compact intervals with the argument
interval above the pole interval.
\end{theorem}

\begin{proof}
Put $P(t)=\prod_j(t-\lambda_j)$, and let $L_i(t)$ be the Lagrange
polynomial of degree $H-1$ with $L_i(d_k)=\mathbf1_{i=k}$.
For any residues $z_h$, the numerator of
$\sum_hz_h/(t-\lambda_{\pi(h)})$ has degree at most $H-1$.
Polynomial interpolation therefore gives
\[
 \sum_h\frac{z_h}{D-\lambda_{\pi(h)}}
 =\sum_i t_i\sum_h\frac{z_h}{d_i-\lambda_{\pi(h)}},
 \qquad t_i=\frac{P(d_i)L_i(D)}{P(D)}. \tag{3}
\]
Every $t_i$ is nonzero by the stated distinctness conditions.
The Cauchy matrix $M_{ij}=1/(d_i-\lambda_j)$ is invertible:
a linear combination of its columns vanishing at all $d_i$ has a
numerator of degree at most $H-1$ with $H$ distinct roots, hence is
zero; its residues at the distinct poles then vanish.
Consequently, for each $h$ there are unique $a_{ih}$ solving
\[
 K_{hj}=\sum_i\frac{t_i a_{ih}}{d_i-\lambda_j}
 \quad\hbox{for all }j.
\]
The matrix in (1) is a column permutation of $M$, so (1) has a unique
solution. Substitute (1) into (3) to obtain (2). Nothing in the argument
requires a lower bound on pole spacing or interval separation beyond
strict separation and distinctness.
\end{proof}

\section{One common height does not read superlinearly many bits}

The preceding theorem makes an arbitrary assignment cost available,
but that cost is shared by all additional rows using the same argument.
This distinction gives an exact obstruction.

\begin{theorem}[Fixed-probe bound]
Let $v_1,\ldots,v_H$ be fixed real numbers. Let $B$ be any function
on the permutations of $H$ symbols; no assignment-cost restriction
on $B$ is needed. For $1\leq q\leq m$, let
\[
 F_q(\pi)=b_q+\sum_h A_{qh}v_{\pi(h)}-B(\pi).
\]
If the strict sign vectors $(\operatorname{sgn}F_q(\pi))_{q=1}^m$
include every element of $\{-1,1\}^m$, then $m\leq H$.
In particular, the asymmetric-anchor construction at one common
extra argument cannot read more than $H$ independent strict bits.
\end{theorem}

\begin{proof}
The vectors
\[
 (v_{\pi(1)},\ldots,v_{\pi(H)},B(\pi))
\]
lie in the affine hyperplane $\sum_hx_h=\sum_jv_j$, whose dimension
is $H$. Each $F_q$ is an affine function on that hyperplane.
More generally, $m$ affine functions $\ell_q(x)=a_q\cdot x+b_q$
on a $d$-dimensional affine space cannot realize all $2^m$ strict
signs if $m>d$. Indeed, choose a nonzero dependence
$\sum_q\alpha_q a_q=0$ and put $k=\sum_q\alpha_q b_q$.
If $k\geq0$, request signs opposite to $\alpha_q$ whenever
$\alpha_q\ne0$; their realization would give
$k=\sum_q\alpha_q\ell_q(x)<0$. If $k<0$, request signs equal to
$\alpha_q$ instead, forcing the same sum to be positive. Assign
arbitrary signs where $\alpha_q=0$. Both cases are contradictions.
Apply this fact with $d=H$.
\end{proof}

\begin{proposition}[A bounded number of common probes]
Suppose the extra rows use arguments from a set of $T$ distinct
values, and for each argument the rows have the preceding form with
one arbitrary shared height and a permutation of one fixed reciprocal
value vector. If every $m$-bit strict sign word occurs, then
$m\leq TH$.
\end{proposition}

\begin{proof}
Concatenate the $T$ lifted value vectors. Each block lies in an affine
space of dimension at most $H$, so their concatenation lies in an
affine space of dimension at most $TH$. Every query is affine on
this space. Apply the dependence argument above. No independence
between the heights or between the blocks is assumed.
\end{proof}

\paragraph{Exact remaining distinction.}
Allowing a growing number of distinct extra arguments evades this
dimension bound. It does not itself supply a decoder. The anchor
contribution at a variable argument is constrained by (3), so it is
not a freely chosen assignment cost at every extra row. A proof of
superlinear decoding must exploit those simultaneous constraints,
rather than replace each row's common height by an independent cost
matrix. This fixed-probe route is therefore blocked as an asymptotic
mechanism, even though its arbitrary-height algebra is valid.

\end{document}
